\clearpage\thispagestyle{empty}
\pdfbookmark[0]{Sources for Observation of the Lunar Eclipse of 31 May 1844 at Göttingen}{lunareclipse1844-sources}
\begin{center}\textbf{Sources for Observation of the Lunar Eclipse of 31 May 1844 at Göttingen}\end{center}
\noindent This AI translation follows the report on p.~466 of \href{https://gdz.sub.uni-goettingen.de/id/PPN236007467}{\textit{Werke}, vol.~VI} (Göttingen, 1874). Carl Friedrich Gauss and Goldschmidt made the observations, and Gauss signed the report. It appeared in \textit{Astronomische Nachrichten}, vol.~XXII, no.~506, \href{https://api.digitale-sammlungen.de/iiif/image/v2/bsb10049447_00024/full/1800,/0/default.jpg}{cols.~31--32}, issued on 20 June 1844. The Bavarian State Library catalogue dates the bound volume to 1845.
\par\smallskip\noindent In the Aristarch ingress row, the collected edition gives Goldschmidt’s minutes as 32, where the journal gives 52. The journal also describes the limb of the Moon farthest from the centre of the shadow, where the collected text says the limb away from the centre. The translation follows the collected text and places Goldschmidt’s 15~1~21 in the Plinius~I row.
\par\smallskip\noindent The short stroke before Goldschmidt’s Dionysius ingress entry is represented as a minus. The first Gauss egress time retains the point after the hour marker, and the Grimaldi~II egress leader retains its comma. Translation and typesetting produced by AI.
\clearpage
